\documentclass[a4paper]{article}

\usepackage[fontset=fandol]{ctex}
\usepackage{amsmath, amssymb}
\usepackage{graphicx}
\usepackage[margin=2.45cm]{geometry}
\usepackage[most]{tcolorbox}
\usepackage{hyperref}
\usepackage{booktabs}
\usepackage{float}
\usepackage{array}
\usepackage{xcolor}
\usepackage{enumitem}
\usepackage{caption}

\setlength{\parindent}{2em}
\setlength{\parskip}{0.35em}
\setlength{\emergencystretch}{3em}
\linespread{1.06}
\sloppy

\newcolumntype{L}[1]{>{\raggedright\arraybackslash}p{#1}}
\newcolumntype{C}[1]{>{\centering\arraybackslash}p{#1}}

\hypersetup{
    colorlinks=true,
    linkcolor=blue!60!black,
    urlcolor=blue!60!black
}

\setlist[itemize]{leftmargin=2.2em,itemsep=0.22em,topsep=0.25em}
\setlist[enumerate]{leftmargin=2.4em,itemsep=0.22em,topsep=0.25em}

\newtcolorbox{knowledgebox}[1]{
    enhanced, colback=blue!5!white, colframe=blue!75!black, colbacktitle=blue!75!black,
    coltitle=white, fonttitle=\bfseries, title={#1},
    attach boxed title to top left={yshift=-2mm, xshift=2mm},
    boxrule=1pt, sharp corners
}

\newtcolorbox{importantbox}[1]{
    enhanced, colback=yellow!10!white, colframe=yellow!80!black, colbacktitle=yellow!80!black,
    coltitle=black, fonttitle=\bfseries, title={#1}, sharp corners
}

\newtcolorbox{warningbox}[1]{
    enhanced, colback=red!5!white, colframe=red!75!black, colbacktitle=red!75!black,
    coltitle=white, fonttitle=\bfseries, title={#1}, sharp corners
}

\newcommand{\notetitle}{异常检测（三）：评估、阈值与代价}
\newcommand{\notesubtitle}{Evaluation, Threshold Selection, and Cost-sensitive Decisions}
\newcommand{\noteauthors}{基于李宏毅老师《机器学习 2025》课程整理}
\newcommand{\notedate}{\today}
\newcommand{\videochannel}{李宏毅深度学习课堂（Bilibili）}
\newcommand{\videoduration}{约 14 分 04 秒}
\newcommand{\videourl}{https://www.bilibili.com/video/BV1TAtwzTE1S?p=31}

\newcommand{\framefig}[4]{%
\begin{figure}[H]
    \centering
    \includegraphics[width=#1\textwidth]{#2}
    \caption{#3\protect\footnotemark}
\end{figure}
\footnotetext{视频画面时间区间：#4。}
}

\begin{document}
\pagestyle{empty}

\begin{center}
    \vspace*{1.1cm}
    {\Huge\bfseries \notetitle\par}
    \vspace{0.35cm}
    {\LARGE \notesubtitle\par}
    \vspace{0.75cm}
    \includegraphics[width=0.62\textwidth]{video.jpg}\par
    \vspace{0.75cm}
    {\large \noteauthors\par}
    {\large \notedate\par}
    \vspace{0.75cm}
    \begin{tcolorbox}[width=0.86\textwidth,center,sharp corners,
                      colback=black!3!white,colframe=black!50!white]
        \centering
        \textbf{视频信息}\\[3pt]
        频道：\videochannel \\
        时长：\videoduration \\
        链接：\href{\videourl}{\videourl}
    \end{tcolorbox}
\end{center}

\newpage
\tableofcontents
\newpage
\pagestyle{plain}

% =====================================================================
\section{从 confidence baseline 到完整评估流程}
% =====================================================================

上一讲已经说明：在有标注的异常检测设定中，可以先训练一个普通分类器，再把分类器给出的 confidence score 当成异常检测分数。本讲继续往下问一个更实际的问题：既然系统已经可以输出 normal 或 anomaly，那么我们如何判断它做得好不好？又该如何决定阈值 $\lambda$？

课程仍用辛普森家族人物分类作例子。训练资料是辛普森家族角色图片，每一张图片 $x$ 都有角色标签 $\hat{y}$，例如它到底是 Homer、Marge、Bart，还是其他辛普森角色。分类器训练完成后，可以对输入图片输出类别预测与 confidence score $c(x)$。异常检测规则写成
\[
f(x)=
\begin{cases}
\mathrm{normal}, & c(x)>\lambda,\\
\mathrm{anomaly}, & c(x)\le \lambda.
\end{cases}
\]
其中 $\lambda$ 是阈值。阈值高，系统会更谨慎，更容易把样本拒绝为异常；阈值低，系统更宽松，更容易接受样本为正常。

\framefig{0.88}{figures/fig01_framework_dev_set.jpg}
{异常检测框架中的 development set：它不再标角色身份，而是标注输入是否来自辛普森家族}
{00:01:45--00:02:00}

\subsection{训练标签与评估标签不一样}

这个例子最容易混淆的地方在于：训练分类器时需要的标签，和评估异常检测器时需要的标签，并不是同一种标签。

训练分类器时，资料全部来自辛普森家族人物。模型需要知道每张图是哪一个角色，所以训练标签是多类别标签：
\[
y \in \{\text{Homer}, \text{Marge}, \text{Bart}, \ldots\}.
\]
但到了异常检测阶段，我们真正关心的不是“它是哪一个辛普森角色”，而是“它是不是辛普森家族人物”。因此 development set 的标签应该是二元标签：
\[
z \in \{\text{from Simpsons}, \text{not from Simpsons}\}.
\]

这件事很重要。假如 development set 仍然只包含辛普森角色身份标签，我们最多只能评估分类器在 closed-set classification 上分得准不准，却无法评估它遇到非辛普森图片时会不会拒绝。异常检测要模拟的是测试时可能出现的开放世界：输入可能来自训练分布，也可能来自训练分布之外。

\begin{importantbox}{development set 的角色}
    Development set 不是拿来重新训练分类器的主要资料，而是拿来调异常检测系统。它应该尽量模拟未来测试时的情境：里面既要有正常样本，也要有异常样本；标签要回答“是否来自正常分布”，而不是只回答“正常分布内部是哪一类”。
\end{importantbox}

\subsection{从 development set 到 testing set}

有了 development set 之后，我们可以把完整系统 $f(x)$ 放到这个集合上计算表现。这里的完整系统包括分类器、confidence score 定义、阈值 $\lambda$，以及任何额外的超参数。换句话说，development set 看的不是单独的分类准确率，而是整个异常检测流程在模拟测试环境中的效果。

\framefig{0.88}{figures/fig02_framework_testing_set.jpg}
{用 development set 计算系统表现，并据此选择阈值 $\lambda$ 与其他超参数；最终再放到 testing set 上评估}
{00:03:00--00:03:15}

流程可以整理成四步：
\begin{enumerate}
    \item 用训练集训练一个分类器，使其能够区分辛普森家族内部的角色。
    \item 从分类器输出中定义 confidence score，例如最大 softmax 分数或其他 confidence 估计。
    \item 在 development set 上调阈值 $\lambda$，并比较不同系统设定的表现。
    \item 固定所有选择之后，在 testing set 或真实线上资料中评估最终表现。
\end{enumerate}

这个流程背后的原则和普通机器学习相同：不能一边看测试集结果一边改阈值，否则测试集就变成了 development set，最终数字会过于乐观。真正可信的测试表现，应该来自所有设计选择都固定之后的独立资料。

\subsection{本章小结}

异常检测系统不是只要会输出一个 confidence score 就结束了。训练集负责学会正常分布内部的类别结构，development set 负责模拟“正常/异常”二元判断并选择阈值，testing set 才负责评估最终系统。特别要记住：异常检测的 development label 应该回答“是否来自正常分布”，而不只是回答正常类别内部的身份。

% =====================================================================
\section{观察 confidence 分布：阈值为什么难选}
% =====================================================================

接下来课程用一个具体 development set 说明阈值选择为什么不只是形式问题。这个集合中有 100 张辛普森家族图片，也有 5 张不是辛普森家族的异常图片。分类器对每张图片都会给出 confidence score，系统再根据阈值决定 normal 或 anomaly。

\framefig{0.88}{figures/fig03_evaluation_confidence_distribution.jpg}
{100 张正常图片与 5 张异常图片的 confidence 分布：大多数正常样本分数高于 0.998，但异常样本也可能拿到不低的分数}
{00:06:15--00:06:30}

\subsection{理想情形与现实情形}

如果 confidence score 完美，那么所有正常样本都会得到很高分，所有异常样本都会得到很低分。此时只要在两群分数之间放一个阈值，就能把 normal 和 anomaly 完全分开。但真实情形通常没有这么干净。

图中大多数辛普森家族图片的 confidence 都非常高，许多甚至高于 $0.998$。这说明分类器对训练分布内的角色相当有把握。可是，非辛普森图片并不一定都会得到低分。有些异常图片仍然可能被分类器看成某个已知角色的脸、颜色或局部形状，从而得到较高 confidence。

因此，confidence score 更像是一条有用但有噪声的线索，而不是一条绝对可靠的边界。阈值 $\lambda$ 的作用，就是把这条连续分数转换成二元决策。转换之后，分数分布中任何重叠区域都会变成错误。

\subsection{阈值移动带来的偏好变化}

阈值 $\lambda$ 越高，系统越严格：只有非常高 confidence 的样本才会被接受为 normal。这样做通常会减少漏掉异常的机会，但也会让部分正常样本被误判为 anomaly。阈值 $\lambda$ 越低，系统越宽松：更多样本会被接受为 normal。这样做会减少正常样本被打扰的情况，但异常样本更容易混进来。

可以把阈值理解成系统的“保守程度”旋钮：
\begin{center}
\begin{tabular}{L{3.0cm} L{4.8cm} L{4.8cm}}
\toprule
阈值设定 & 倾向 & 常见后果 \\
\midrule
较高的 $\lambda$ & 更容易拒绝输入 & false alarm 增加，missing 可能减少 \\
较低的 $\lambda$ & 更容易接受输入 & false alarm 减少，missing 可能增加 \\
随任务调节的 $\lambda$ & 按应用代价折中 & 需要 development set 或业务代价信息 \\
\bottomrule
\end{tabular}
\end{center}

\begin{knowledgebox}{不要把分数大小直接等同于答案}
    Confidence score 是排序与决策的依据，但它本身不是最终答案。只有当阈值、错误代价和应用目标一起确定后，我们才知道某个分数到底应该被解释成 normal 还是 anomaly。
\end{knowledgebox}

\subsection{本章小结}

在实际资料中，正常样本和异常样本的 confidence 分布往往会重叠。大多数正常样本分数高，并不代表所有异常样本分数都低。异常检测的困难正在于此：阈值不是从图上随便画一条线，而是在 false alarm 与 missing 之间选择一种任务偏好。

% =====================================================================
\section{Accuracy 的陷阱与两类错误}
% =====================================================================

普通分类任务常用 accuracy 衡量系统表现，但课程特别强调：accuracy 不是异常检测中的好指标。原因并不是 accuracy 的定义错了，而是异常检测资料通常高度不平衡。正常资料很多，异常资料很少；在这种设定下，一个几乎不做事的系统也可能拿到很高 accuracy。

\framefig{0.88}{figures/fig04_accuracy_not_good_measurement.jpg}
{Accuracy 可能误导异常检测评估：系统把 5 个异常全放过，仍可因 100 个正常样本正确而得到 95.2\%}
{00:08:15--00:08:30}

\subsection{高 accuracy 可能代表系统什么都没做}

课程中的例子有 100 张正常图片和 5 张异常图片。假设某个系统把所有样本都判断成 normal，那么 100 张正常图片都判断对了，5 张异常图片都判断错了。它的 accuracy 是
\[
\frac{100}{100+5}=95.2\%.
\]
这个数字看起来很高，但系统完全没有侦测出任何异常。若用户真正关心的是“异常出现时能不能抓到”，那么这样的系统几乎没有价值。

这就是异常检测评估的核心陷阱：当 anomaly 很稀少时，accuracy 会被正常样本数量主导。只要正常样本够多，模型即使漏掉所有异常，也可能得到一个漂亮的总准确率。

\begin{warningbox}{类别不平衡会掩盖失败}
    异常检测最怕的不是总错误率看起来高，而是关键错误被大量正常样本稀释。一个系统如果永远回答 normal，在许多极端不平衡任务中都能得到高 accuracy，但它没有完成异常检测的任务。
\end{warningbox}

\subsection{两种错误：false alarm 与 missing}

为了更细致地描述异常检测表现，课程区分两类错误：
\begin{itemize}
    \item \textbf{False alarm}：正常资料被判断成异常。也就是系统误报，可能造成额外检查、人工复核或用户困扰。
    \item \textbf{Missing}：异常资料被判断成正常。也就是系统漏报，可能让真正危险或重要的异常事件没有被处理。
\end{itemize}

\framefig{0.88}{figures/fig05_false_alarm_and_missing.jpg}
{同一个阈值会同时决定 false alarm 与 missing：正常样本被拒绝是误报，异常样本被接受是漏报}
{00:09:30--00:09:45}

这两类错误没有绝对的谁更严重，严重程度取决于任务。如果是垃圾邮件过滤，误把重要邮件判成垃圾邮件可能非常烦；如果是安检或医疗筛查，漏掉危险案例可能更严重。异常检测评估不能只问“错了几个”，还要问“错的是哪一种”。

\subsection{用混淆矩阵看错误结构}

异常检测的二元决策可以写成混淆矩阵。这里用“Detected”表示系统判断为 anomaly，用“Not Det”表示系统判断为 normal：
\begin{center}
\begin{tabular}{L{3.0cm} C{3.0cm} C{3.0cm} L{4.2cm}}
\toprule
系统输出 & 真实 anomaly & 真实 normal & 含义 \\
\midrule
Detected & 正确侦测 & false alarm & 把输入当成异常 \\
Not Det & missing & 正确接受 & 把输入当成正常 \\
\bottomrule
\end{tabular}
\end{center}

相比单一 accuracy，混淆矩阵能告诉我们错误是如何分布的：系统到底是太敏感，还是太宽松；问题主要是误报太多，还是漏报太多。后面讨论 cost table 时，正是基于这个结构计算不同错误的代价。

\subsection{本章小结}

异常检测不能只看 accuracy，因为异常样本通常太少，导致“全部判断为 normal”的系统也可能得到高分。更有用的评估方式是把错误拆成 false alarm 和 missing：前者是正常样本被误报成异常，后者是异常样本被漏掉成正常。不同应用对两种错误的容忍度不同，因此需要更贴近任务的指标。

% =====================================================================
\section{用 cost table 决定系统偏好}
% =====================================================================

既然 false alarm 和 missing 的重要性会随任务改变，课程进一步引入 cost table。Cost table 的想法很直接：不要假设每一种错误都一样糟，而是给不同预测结果指定不同代价，再计算系统的总成本。总成本较低的系统，才是当前任务偏好下较好的系统。

\framefig{0.88}{figures/fig06_two_thresholds_confusion_matrices.jpg}
{两个阈值产生两组混淆矩阵：左边较少 false alarm 但漏掉更多异常，右边抓到更多异常但误报也更多}
{00:10:45--00:11:00}

\subsection{两个阈值，对应两种系统行为}

课程展示了两个阈值选择。左边系统的混淆矩阵为：
\[
\begin{array}{c|cc}
 & \text{真实 anomaly} & \text{真实 normal} \\
\hline
\text{Detected} & 1 & 1\\
\text{Not Det} & 4 & 99
\end{array}
\]
它只抓到 1 个异常，漏掉 4 个异常；同时只有 1 个正常样本被误报为异常。这个系统比较保守，不太打扰正常样本，但会漏掉较多异常。

右边系统的混淆矩阵为：
\[
\begin{array}{c|cc}
 & \text{真实 anomaly} & \text{真实 normal} \\
\hline
\text{Detected} & 2 & 6\\
\text{Not Det} & 3 & 94
\end{array}
\]
它抓到 2 个异常，漏掉 3 个异常；但有 6 个正常样本被误报。这个系统比较敏感，愿意付出更多 false alarm 来换取较少 missing。

\subsection{Cost Table A：误报很贵}

第一种 cost table 假设 false alarm 代价很高：把正常样本误判为异常要付出 100 的成本，而漏掉异常只付出 1 的成本。可以写成：
\[
\begin{array}{c|cc}
 & \text{真实 anomaly} & \text{真实 normal} \\
\hline
\text{Detected} & 0 & 100\\
\text{Not Det} & 1 & 0
\end{array}
\]
于是左边系统的总成本是
\[
1\times 100 + 4\times 1 = 104,
\]
右边系统的总成本是
\[
6\times 100 + 3\times 1 = 603.
\]
在这个代价设定下，左边系统明显更好，因为它大幅减少了误报。

\subsection{Cost Table B：漏报很贵}

第二种 cost table 反过来，假设 missing 的代价很高：漏掉异常要付出 100 的成本，而 false alarm 只付出 1 的成本：
\[
\begin{array}{c|cc}
 & \text{真实 anomaly} & \text{真实 normal} \\
\hline
\text{Detected} & 0 & 1\\
\text{Not Det} & 100 & 0
\end{array}
\]
此时左边系统的总成本变成
\[
1\times 1 + 4\times 100 = 401,
\]
右边系统的总成本变成
\[
6\times 1 + 3\times 100 = 306.
\]
结论完全反过来：右边系统更好。它虽然误报更多，但少漏掉一个异常；在 missing 很贵的任务里，这个交换是值得的。

\framefig{0.88}{figures/fig07_cost_tables_change_preference.jpg}
{同两套系统在 Cost Table A 与 Cost Table B 下优劣相反：评价标准会改变我们选择的模型或阈值}
{00:12:15--00:12:30}

\subsection{医学筛查中的直觉}

课程用癌症检测解释为什么 cost table 不能脱离任务。假设系统判断一张检查影像是否异常。如果病人其实没有癌症，但系统误判为有癌症，这是 false alarm。它可能造成进一步检查、心理压力和医疗成本，但通常还有后续流程可以纠正。

反过来，如果病人其实有癌症，系统却判断为正常，这就是 missing。它可能导致治疗延迟，后果远比一次额外检查严重。因此在癌症筛查这类任务里，missing 通常比 false alarm 更贵。系统宁愿多发出一些警报，也不愿轻易放过真正异常。

\begin{importantbox}{指标不是中立的}
    选择评价指标就是选择任务偏好。Accuracy 假设每个样本、每种错误都同等重要；cost table 则明确告诉模型评估过程：哪些错误更不能接受。异常检测往往需要后者。
\end{importantbox}

\subsection{本章小结}

Cost table 把异常检测从“数错几个”推进到“错在哪里、代价多大”。同样两个系统，在 false alarm 昂贵的任务里可能选择左边，在 missing 昂贵的任务里可能选择右边。因此，阈值和模型不能脱离应用场景讨论；评估标准本身就是系统设计的一部分。

% =====================================================================
\section{不固定阈值的评价：ranking 与 AUC}
% =====================================================================

前面的讨论都假设系统最终要选择一个阈值 $\lambda$，再把样本分成 normal 和 anomaly。但有些评价指标并不先固定阈值，而是只关心 anomaly score 的排序。课程最后提到的典型例子是 ROC 曲线下面积，也就是 AUC。

\framefig{0.88}{figures/fig08_auc_ranking_metric.jpg}
{某些评价指标只考虑异常分数的排序，例如 ROC 曲线下面积 AUC，因此不需要预先固定一个阈值}
{00:13:45--00:14:00}

\subsection{从二元决策到排序}

如果系统输出的是 confidence score $c(x)$，通常 confidence 越低越可能是异常。为了做排序，可以定义 anomaly score：
\[
s(x)=-c(x)
\]
或使用任何“分数越高越异常”的等价形式。Ranking metric 关心的是：真正的异常样本是否整体排在正常样本前面。它不急着问“阈值放在哪里”，而是问“如果从最可疑到最不可疑排序，异常会不会被排到前面”。

这种观点很适合还没有确定部署阈值的阶段。研究者可以先比较不同模型产生的排序质量，再根据实际应用选择工作点。例如医疗筛查可以选择高召回的阈值，内容审核可以选择人工审核量可承受的阈值，设备监控可以选择报警频率不过载的阈值。

\subsection{ROC 与 AUC 的含义}

ROC 曲线会扫描所有可能阈值，观察不同阈值下 true positive rate 与 false positive rate 的关系。AUC 是 ROC 曲线下面积，可以理解为一种整体排序能力的摘要。直观地说，AUC 越高，表示系统越倾向于把真实 anomaly 排在真实 normal 前面。

更形式化一点，AUC 可以解释成：
\[
\Pr\bigl(s(x_{\text{anomaly}}) > s(x_{\text{normal}})\bigr),
\]
也就是随机抽一个异常样本和一个正常样本时，系统把异常样本打得更可疑的概率。这个解释说明了为什么 AUC 不需要固定阈值：它评价的是成对排序，而不是某一个阈值下的二元分类。

\subsection{AUC 能回答什么，不能回答什么}

AUC 的优点是避免提前选择 $\lambda$，并且比 accuracy 更不容易被单一类别数量支配。它特别适合模型开发早期，用来判断一个异常分数是否有区分能力。

但 AUC 也不是部署答案。真实系统最终往往仍要决定报警阈值、人工复核量、可接受漏报率，以及不同错误的业务代价。两个模型 AUC 接近时，某个模型可能在高召回区域更好，另一个模型可能在低误报区域更好。若应用只关心某个特定工作区间，完整 ROC 曲线、precision-recall 曲线或 cost-sensitive evaluation 会比单一 AUC 更有信息。

\begin{knowledgebox}{开发阶段与部署阶段的指标}
    开发阶段可以用 AUC 评价排序能力，因为它不依赖某个固定阈值。部署阶段通常还需要根据业务代价选择具体阈值，并报告该阈值下的 false alarm、missing、召回率或成本。
\end{knowledgebox}

\subsection{本章小结}

AUC 代表一种不先固定阈值的评价思路：先看模型是否能把异常样本排到正常样本前面。它适合比较 anomaly score 的整体质量，但不能替代最终部署时的阈值选择和代价分析。异常检测评估通常需要同时考虑排序能力、工作点表现和应用代价。

% =====================================================================
\section{总结与延伸}
% =====================================================================

这一讲的主线是：异常检测的困难不只在模型本身，也在评估方式。一个分类器可以给出 confidence score，一个阈值可以把分数变成 normal 或 anomaly，但系统是否“好”，取决于我们如何定义错误、如何看待类别不平衡，以及如何衡量不同错误的代价。

可以把本讲内容压缩成四个判断：
\begin{enumerate}
    \item Development set 必须模拟 anomaly detection 的测试情境，标签应是“是否来自正常分布”，不是正常类别内部身份。
    \item Accuracy 在异常检测中经常误导，因为正常样本数量可能远大于异常样本数量。
    \item False alarm 与 missing 是两种不同错误，哪个更严重由应用场景决定。
    \item AUC 等 ranking metric 可以在不固定阈值时评价分数质量，但部署时仍要选择具体工作点。
\end{enumerate}

如果把这套逻辑迁移到真实项目，一个较稳妥的异常检测报告不应该只写“accuracy = 95\%”。它至少应该说明：测试集中正常与异常比例是多少；阈值如何选择；false alarm 和 missing 分别有多少；如果任务中两种错误代价不同，总成本如何计算；如果还没有固定阈值，则应该报告 AUC 或其他排序指标。

从下一步学习角度看，本讲也埋下了一个更一般的问题：异常检测不是单纯追求更复杂的模型，而是要把模型分数、资料分布、阈值选择和应用代价连在一起。只有评估方式对了，模型改进才有方向；否则模型可能只是把一个漂亮但无用的数字变得更漂亮。

\end{document}
